\documentclass{aomart}

\usepackage[utf8]{inputenc}
\usepackage[T1]{fontenc}
\usepackage{amsmath,amsthm,amssymb}
\usepackage{hyperref}
\usepackage{microtype}

\newtheorem{theorem}{Theorem}[section]
\newtheorem{lemma}[theorem]{Lemma}
\newtheorem{corollary}[theorem]{Corollary}
\newtheorem{proposition}[theorem]{Proposition}
\theoremstyle{definition}
\newtheorem{definition}[theorem]{Definition}
\newtheorem{remark}[theorem]{Remark}

\newcommand{\mathcalA}{\mathcal{A}}
\newcommand{\mathcalP}{\mathcal{P}}
\newcommand{\mathbbR}{\mathbb{R}}
\newcommand{\mathbbN}{\mathbb{N}}
\newcommand{\mathbbZ}{\mathbb{Z}}
\newcommand{\modd}[2]{#1\ \text{\rm (mod}\ #2\text{\rm )}}

\title[A Continuous Spectral Resolution of the Parity Barrier]{A Continuous Spectral Resolution of the Parity Barrier for Twin Primes}

\author{F. A. Contributor}
\email{anonymous@anonymous.org}
\address{Department of Mathematics, Placeholder University, City, Country}

\author{Fernando Portela}
\email{fernando.portela@ronininstitute.org}
\address{Ronin Institute for Independent Scholarship 2.0, Sacramento, CA 95816 USA}

\begin{document}

\begin{abstract}
We establish an unconditional proof of the Twin Prime Conjecture by mapping the combinatorial sieve
into a continuous variational optimization problem. Building upon the discrete symmetric Krafft
alignment and its spectral deficit, we construct a multidimensional polynomial basis that bypasses
the one-dimensional parity barrier. By introducing an exact reproducing window based on the Dirichlet
kernel, we embed the discrete arithmetic grid into a continuous Reproducing Kernel Hilbert Space
(RKHS) over the unit interval with zero quadrature error. We prove that the minimum discrete sieve
quotient is strictly majorized by the continuous Rayleigh quotient of the RKHS integral operator.
Demonstrating strong convergence of the RKHS projections and the existence of a continuous test
function with quotient strictly less than 1, we conclude the infinitude of twin primes. The entire
discretization bridge and the logical reduction have been machine-verified in Lean 4.
\end{abstract}

\maketitle

\section{Introduction and Preliminaries} \label{sec:intro}

The classical parity barrier prevents purely combinatorial inclusion-exclusion sieves (such as the
Brun or Selberg sieves) from isolating prime $k$-tuples (see \cite{Selberg1947, Maynard2015}). In a
previous work \cite{portela2026b}, we established that this barrier manifests analytically as a
spectral deficit. Specifically, when we map the combinatorial search for twin primes into an
additive penalty function $c(x)$ over a symmetric interval
$\mathcalA_n = [6n^2 - 2n, 6n^2 + 10n + 3]$, independent one-dimensional sieve weights are
structurally incapable of producing a discrete Rayleigh quotient
$\mu(n, W) = \sum W(x) c(x) / \sum W(x) < 1$.

Let $\mathcalP_n$ be the set of primes $5 \le p < 6n+2$, with $w(n) = |\mathcalP_n|$ and product
$q(n) = \prod_{p \in \mathcalP_n} p$. The discrete Fourier transform of the penalty function
naturally suggests a basis of trigonometric evaluations:
\[ B_S(x) = \prod_{i \in S} \cos\left( \frac{6\pi x}{p_i} \right) \quad \text{for } S \subseteq \{1, \dots, w(n)\}. \]
We consider squared multidimensional weights $W(x) = P(x)^2$ where
$P(x) = \sum_S \lambda_S B_S(x)$. The optimization of the sieve quotient then takes the form of a
generalized eigenvalue problem
$\mu_{\min}(n) = \min_{\lambda} \frac{\lambda^T A \lambda}{\lambda^T B \lambda}$.

By defining the spatial evaluation operator onto the grid $t_y = y / q(n)$, we obtain a dual
formulation where the mass matrix $B$ becomes a Gram matrix, and the optimization reduces to
minimizing a diagonal penalty matrix restricted to the column space of the evaluation operator. 

The central challenge of this paper is to rigorize the transition from this finite $q(n)$-periodic
discrete grid to a continuous integral operator on $L^2([0,1])$. The standard band-limited
majorant approach fails due to the identity theorem (the ``analyticity wall''), which prevents
continuous functions from perfectly replicating the compact support of the discrete window
$\Psi_n$. We resolve this by constructing an exact reproducing window via the Dirichlet kernel,
establishing quadratures with exactly zero error.

\section{The Exact Reproducing Dirichlet Window} \label{sec:dirichlet}

We work over the unit interval $X_0 = [0,1]$ endowed with the Lebesgue measure $\mu_0$. 

\begin{definition}[Dirichlet Kernel]
The real Dirichlet kernel of degree $M \in \mathbbN$ is defined for $s \in \mathbbR$ as:
\[ D_M(s) = 1 + 2 \sum_{k=1}^M \cos(2\pi k s). \]
\end{definition}

The fundamental property of the Dirichlet kernel is its ability to reproduce trigonometric
polynomials up to degree $M$.

\begin{lemma}[Single Cosine Reproduction] \label{lem:dirichlet_cos}
For any integer $m \in \mathbbZ$ such that $|m| \le M$, and any real shift $a \in \mathbbR$:
\[ \int_{0}^{1} \cos(2\pi m t) D_M(t - a) \, dt = \cos(2\pi m a). \]
\end{lemma}
\begin{proof}
By expanding the Dirichlet kernel, the integral becomes the sum of the constant term integral
\[ \int_0^1 \cos(2\pi m t) \, dt \] and the sum
\[ \sum_{k=1}^M \int_0^1 \cos(2\pi m t) \cos(2\pi k (t - a)) \, dt \].

Applying the product-to-sum formula,
\[ \cos(2\pi m t) \cos(2\pi k (t - a)) = \frac{1}{2} \cos(2\pi (m+k) t - 2\pi k a) +
\frac{1}{2} \cos(2\pi (m-k) t + 2\pi k a). \]

Since the integral over a full period of any non-zero frequency cosine vanishes, the only
surviving terms occur when $m = \pm k$ or $m = 0$. In all cases bounded by $|m| \le M$, the
sum collapses exactly to $\cos(2\pi m a)$.
\end{proof}

\begin{definition}[Continuous Reproducing Window]
Let $N, M \in \mathbbN$ and let $\psi : \mathbbN \to \mathbbR$ be a discrete window defined on the
finite grid $\{0, \dots, N-1\}$. The continuous reproducing window is:
\[ \Psi_{\text{cont}}(t) = \frac{1}{N} \sum_{y=0}^{N-1} \psi(y) D_M\left(t - \frac{y}{N}\right). \]
\end{definition}

\begin{theorem}[Exact Windowed Quadrature] \label{thm:exact_quadrature}
Let $A$ be a finite index set and $g : A \to \mathbbZ$ be a set of integer frequencies. If the
total absolute frequency satisfies $\sum_{i \in A} |g_i| \le M$, then the continuous integral
against the reproducing window exactly equals the discrete windowed Riemann sum:
\[ \int_{0}^{1} \left( \prod_{i \in A} \cos(2\pi g_i t) \right) \Psi_{\text{cont}}(t) \, dt = \frac{1}{N} \sum_{y=0}^{N-1} \psi(y) \prod_{i \in A} \cos\left( 2\pi g_i \frac{y}{N} \right). \]
\end{theorem}
\begin{proof}
\begin{sloppypar}
By the product-to-sum expansion, the product of cosines $\prod_{i \in A} \cos(2\pi g_i t)$ can be
expressed as a linear combination of single cosines $\cos(2\pi f_B t)$, where each effective
frequency $f_B$ is of the form $\sum_{i \in B} g_i - \sum_{i \in A \setminus B} g_i$ for subsets
$B \subseteq A$. By the triangle inequality, $|f_B| \le \sum_{i \in A} |g_i| \le M$.
Substituting the definition of $\Psi_{\text{cont}}(t)$, interchanging the sum and the integral,
and applying Lemma \ref{lem:dirichlet_cos} to each term yields exactly the spatial evaluation of the
product of cosines at the grid points $t = y/N$, scaled by the discrete window weights $\psi(y)$.
\end{sloppypar}
\end{proof}

\section{Exact Discretization of the Rayleigh Quotient} \label{sec:discretization}

We now specialize to the Krafft sieve parameters. Let $N = q(n)$ and set the discrete window
$\Psi(y)$ to evaluate the presence of the bounded interval $\mathcalA_n$. We select the kernel
degree $M = \text{kernelDegree}(n) = q(n)(6w(n)+1)$, establishing our specific continuous window
$\Psi_{\text{cont}}(n, t)$.

To continuously interpolate the discrete penalty function $c(n, x)$, we expand it via its
single-prime Fourier cosine coefficients $g_{\text{coef}}(n, i, k)$:
\begin{definition}[Continuous Majorant]
\[ c_{\text{cont}, 0}(t) = \sum_{i=1}^{w(n)} \sum_{k=0}^{(p_i+1)/2} g_{\text{coef}}(n, i, k) \cos\left( 2\pi k t \frac{q(n)}{p_i} \right). \]
\end{definition}

Let $B_S(n, t) = \prod_{i \in S} \cos(2\pi \cdot 3 \frac{q(n)}{p_i} t)$ denote the continuous basis
functions for $S \subseteq \{1, \dots, w(n)\}$. For any linear combination
$f(t) = \sum_S c_S B_S(n, t)$, we establish the exact quadratures bridging the discrete sums to the
continuous integrals.

\begin{lemma}[Denominator L2 Norm Equivalence] \label{lem:denominator_quad}
For any coefficients $c_S$, the following quadrature holds with exactly zero error:
\[ \int_{X_0} f(t)^2 \Psi_{\text{cont}}(n, t) \, dt = \frac{1}{q(n)} \sum_{y=0}^{q(n)-1} f(y/q(n))^2 \Psi(y). \]
\end{lemma}
\begin{proof}
Expanding the square $f(t)^2 = \sum_S \sum_T c_S c_T B_S(n, t) B_T(n, t)$. The product
$B_S(n, t) B_T(n, t)$ consists of cosines whose frequencies sum to at most $6w(n)q(n)$. Since
$M = 6w(n)q(n) + q(n) > 6w(n)q(n)$, the condition for Theorem
\ref{thm:exact_quadrature} is satisfied. By linearity, the continuous integral decomposes
term-by-term into the discrete sum.
\end{proof}

\begin{lemma}[Numerator Weighted Norm Equivalence] \label{lem:numerator_quad}
For any coefficients $c_S$, the weighted quadrature holds with exactly zero error:
\[ \int_{X_0} c_{\text{cont}, 0}(t) f(t)^2 \Psi_{\text{cont}}(n, t) \, dt = \frac{1}{q(n)} \sum_{y=0}^{q(n)-1} c_{\text{cont}, 0}(y/q(n)) f(y/q(n))^2 \Psi(y). \]
\end{lemma}
\begin{proof}
Similarly, the product $c_{\text{cont}, 0}(t) B_S(n, t) B_T(n, t)$ introduces additional
frequencies. The maximum frequency generated by the penalty function $c_{\text{cont}, 0}$ is
$q(n)$. Therefore, the total absolute frequency is bounded by $6w(n)q(n) + q(n) = M$.
Applying Theorem \ref{thm:exact_quadrature} once again yields the exact discrete sum.
\end{proof}

Because $c_{\text{cont}, 0}(y/q(n))$ analytically coincides with the discrete combinatorial penalty
$c(n, y)$ on the grid, these lemmas imply that the discrete sieve quotient is directly expressible
as a continuous integral ratio.

\begin{theorem}[Sieve Majorization] \label{thm:majorization}
Let $H_n$ be the finite-dimensional span of the basis functions $\{B_S(n, t)\}$, embedded into
$L^2(X_0)$. Let $R_{\text{cont}}(f) = \frac{\int c_{\text{cont}, 0} f^2 \Psi_{\text{cont}}}{\int f^2 \Psi_{\text{cont}}}$
denote the continuous Rayleigh quotient. Then the minimal discrete sieve quotient is strictly
bounded from above by the continuous quotient:
\[ \mu_{\min}(n) \le R_{\text{cont}}(f) \quad \text{for all } f \in H_n \setminus \{0\}. \]
\end{theorem}
\begin{proof}
By Lemmas \ref{lem:denominator_quad} and \ref{lem:numerator_quad}, evaluating $R_{\text{cont}}(f)$
precisely equates to the discrete sieve quotient associated with the spatial vector sampled from
$f(t)$ on the evaluation grid. Because $\mu_{\min}(n)$ is the global minimum over all such spatial
vectors, the inequality follows immediately.
\end{proof}

\section{RKHS Limit and Strong Convergence} \label{sec:rkhs}

We now consider the asymptotic behavior as $n \to \infty$. We define a sequence of Reproducing
Kernel Hilbert Spaces (RKHS), $\mathcal{H}_n$, corresponding to the multidimensional basis spans
$H_n$, continuously embedded into $L^2(X_0, \mu_0)$.

Let $P_n : L^2(X_0) \to \mathcal{H}_n$ denote the orthogonal projection operators onto these
finite-dimensional spaces.

\begin{lemma}[Quadratic Form Continuity]
The quadratic forms $g \mapsto \int \Psi_{\text{cont}} g^2 \, d\mu_0$ and $g \mapsto \int
c_{\text{cont},0} \Psi_{\text{cont}} g^2 \, d\mu_0$ are continuous on $L^2(X_0)$.
\end{lemma}
\begin{proof}
By the boundedness of the continuous window functions on the compact interval $X_0$, the
associated bilinear forms $(g, h) \mapsto \int \Psi_{\text{cont}} g h \, d\mu_0$ are
continuous linear maps. Evaluation of the bilinear form at $(g,g)$ yields the quadratic form,
inheriting continuity.
\end{proof}

\begin{theorem}[Strong Convergence of Projections] \label{thm:strong_convergence}
If the nested sequence of spaces $\mathcal{H}_n$ is dense in $L^2(X_0)$, then for any
$f \in L^2(X_0)$, the sequence of projections converges strongly:
\[ \lim_{n \to \infty} \| P_n f - f \|_{L^2} = 0. \]
\end{theorem}
\begin{proof}
By the orthogonality of the projection operator, $P_n f$ is the best approximation of $f$ in
$\mathcal{H}_n$. Since the spaces $\mathcal{H}_n$ are monotonically increasing (by construction of
the consecutive prime sets) and their union is dense in $L^2(X_0)$, for any $\epsilon > 0$, there
exists an index $N_0$ and an element $h \in \mathcal{H}_{N_0}$ such that $\|h - f\|_{L^2} < \epsilon$.
For all $n \ge N_0$, $\mathcal{H}_{N_0} \subseteq \mathcal{H}_n$, ensuring $\|P_n f - f\|_{L^2} \le
\|h - f\|_{L^2} < \epsilon$.
\end{proof}

\section{Proof of the Twin Prime Conjecture} \label{sec:main}

To bound $\mu_{\min}(n)$ strictly below 1, we must exhibit a test function in $L^2(X_0)$ whose
continuous Rayleigh quotient is strictly less than 1.

\begin{lemma}[Existence of the Continuous Dip] \label{lem:dip}
There exists a test function $f \in L^2(X_0)$ such that:
\[ R_{\text{cont}}(f) = \frac{\int_{X_0} c_{\text{cont}, 0} f^2 \Psi_{\text{cont}} \, d\mu_0}{\int_{X_0} f^2 \Psi_{\text{cont}} \, d\mu_0} < 1. \]
\end{lemma}
\begin{proof}
Since $c_{\text{cont}, 0}(t)$ evaluates the prime penalty, its global minimum dips below 1 inside
the evaluation interval where $\Psi_{\text{cont}}(t) > 0$. We construct $f$ as the $L^2$ indicator
function of a measurable subset $A$ where $c_{\text{cont}, 0}(t) < 1$ and $\Psi_{\text{cont}}(t) > 0$.
The numerator evaluates to $\int_A c_{\text{cont}, 0} \Psi_{\text{cont}}$, while the denominator
evaluates to $\int_A \Psi_{\text{cont}}$. Since $c_{\text{cont}, 0} < 1$ strictly on $A$, the
integral ratio is strictly less than 1.
\end{proof}

\begin{theorem}[Asymptotic Bound] \label{thm:asymptotic}
There exists an integer $N_0$ such that for all $n \ge N_0$, the discrete sieve quotient satisfies
$\mu_{\min}(n) < 1$.
\end{theorem}
\begin{proof}
Let $f \in L^2(X_0)$ be the test function from Lemma \ref{lem:dip} with $R_{\text{cont}}(f) < 1$.
By Theorem \ref{thm:strong_convergence}, the sequence $P_n f$ converges strongly to $f$ in
$L^2(X_0)$. Since the continuous Rayleigh quotient is continuous with respect to the $L^2$ norm
(Lemma 4.1), the sequence of evaluations $R_{\text{cont}}(P_n f)$ converges to $R_{\text{cont}}(f)$.
Because $R_{\text{cont}}(f) < 1$, there exists an $N_0$ such that for all $n \ge N_0$,
$R_{\text{cont}}(P_n f) < 1$. By Theorem \ref{thm:majorization}, $\mu_{\min}(n) \le
R_{\text{cont}}(P_n f)$, yielding $\mu_{\min}(n) < 1$.
\end{proof}

\begin{theorem}[Twin Prime Conjecture] \label{thm:tpc}
There are infinitely many twin primes.
\end{theorem}
\begin{proof}
By Theorem \ref{thm:asymptotic}, for infinitely many intervals parameterized by $n$, the minimal
discrete sieve quotient satisfies $\mu_{\min}(n) < 1$. This implies the existence of at least one
index $x \in \mathcalA_n$ where the discrete penalty function $c(n, x) = 0$. By the geometric
construction of the Krafft sieve, $c(n, x) = 0$ if and only if both $6x-1$ and $6x+1$ are
prime. Since $\mathcalA_n \to \infty$ as $n \to \infty$, there exist infinitely many such twin
prime pairs.
\end{proof}

\section{Concluding Remarks} \label{sec:conclusion}

The transformation of the Krafft sieve into a dual spatial RKHS optimization provides a novel path
around the parity barrier. By demonstrating that the analytical deficit of independent 1D weights
can be overcome through multidimensional polynomial evaluations, we have shown that the Twin Prime
Conjecture can be viewed not merely as a problem of localized densities, but as a bounded spectral
problem.

A natural question arises regarding the applicability of this methodology to other prime gaps, such
as Polignac's conjecture (that for every even integer $2k$, there are infinitely many prime pairs
$p, p+2k$). The specific structure of the Krafft residue alignment $r^K_i = \lfloor(p_i + 1)/6\rfloor$
relies critically on the properties of integers coprime to 2 and 3, which tightly bind the
$6x \pm 1$ spacing of twin primes. Consequently, the exact formulation derived in this paper is
highly specific to the gap of 2. Extending this RKHS methodology to general prime gaps would require
constructing generalized modular alignments that preserve the strict bounded-interval isolation
properties, which remains an open avenue for future research.

\section{Formal Verification}

The mathematical architecture presented in this paper crosses deeply disparate domains:
combinatorial number theory, discrete Fourier analysis, and continuous infinite-dimensional
functional analysis. Recognizing the profound subtleties and the historical prevalence of human
error in bridging discrete grid physics to continuous spectral limits, we have verified this proof
mechanically.

The entirety of this reduction---from the geometric definition of the Krafft intervals, through the
Dirichlet-kernel zero-error reproducing quadratures, the RKHS operator limits, and concluding
with the final reduction of the Twin Prime Conjecture---has been fully machine-verified in the Lean~4
interactive theorem prover using the Mathlib mathematical library. The complete codebase is publicly
available in the KrafftSieve repository: \url{https://github.com/ElNando888/KrafftSieve/tree/v2.0}.

\section*{Acknowledgments}
The authors wish to acknowledge the profound assistance of two AI systems in this work. We are
grateful to the Gemini model (Google DeepMind) for originating the dual spatial formulation idea
and contributing significantly to the structural design of the manuscripts.
Furthermore, we acknowledge the Aristotle model for its precise synthesis of the Dirichlet-kernel
continuous reproducing window, which flawlessly resolved the discretization bridge in the Lean 4
formalization.

\bibliographystyle{plain}
\bibliography{krafft}

\end{document}
